\begin{UseCase}{CU-57}{Bloqueo automático de sesión por inactividad}{
	Permite al sistema cerrar automáticamente la sesión de un usuario cuando no ha registrado actividad durante 30 minutos consecutivos, protegiendo la confidencialidad de la información del expediente.
}
	\UCitem{Versión}{\color{Gray}1.0}
	\UCitem{Autor}{\color{Gray}Ramírez Cuevas Emiliano}
	\UCitem{Supervisa}{\color{Gray}Rivera Segura José Emiliano}
	\UCitem{Actor}{Sistema (proceso automático)}
	\UCitem{Propósito}{Proteger la información confidencial de los expedientes ante situaciones de abandono no intencional de la sesión.}
	\UCitem{Entradas}{Tiempo de última actividad del usuario en sesión.}
	\UCitem{Origen}{Proceso automático del servidor de sesión.}
	\UCitem{Salidas}{Sesión expirada y redirección a pantalla de inicio de sesión con mensaje de aviso.}
	\UCitem{Destino}{Pantalla de inicio de sesión.}
	\UCitem{Precondiciones}{El usuario debe tener una sesión activa con inactividad mayor a 30 minutos.}
	\UCitem{Postcondiciones}{La sesión queda destruida y el usuario debe autenticarse nuevamente.}
	\UCitem{Errores}{Ninguno.}
	\UCitem{Tipo}{Caso de uso secundario}
	\UCitem{Observaciones}{Se muestra un mensaje de advertencia 2 minutos antes del cierre automático.}
\end{UseCase}

\begin{UCtrayectoria}
	\UCpaso El sistema detecta que el usuario lleva 28 minutos sin registrar actividad.
	\UCpaso Muestra un mensaje de advertencia \hyperlink{mMSG03}{\bf MSG-03} indicando que la sesión se cerrará en 2 minutos por inactividad.
	\UCpaso[\UCactor] Si el usuario retoma la actividad, el contador se reinicia. De lo contrario, el sistema continúa al paso 4.
	\UCpaso Transcurridos los 30 minutos totales, destruye la sesión del usuario.
	\UCpaso Redirige al usuario a la pantalla de inicio de sesión mostrando el mensaje de sesión expirada.
	\UCpaso[] -- -- Fin del caso de uso.
\end{UCtrayectoria}
